\documentclass[11pt]{article}
\newcommand{\doctitle}{Matching the Analysis Model to the
  Data-Generating Process:
  AR(1) Residual Correlation in the
  pmsimstats LME Analysis}
\newcommand{\shorttitle}{AR(1) Residual Correlation}
\newcommand{\docdate}{2026-03-21}
% pmsimstats-ng standard document preamble
% Usage: % pmsimstats-ng standard document preamble
% Usage: % pmsimstats-ng standard document preamble
% Usage: \input{pmsimstats-preamble} after \documentclass[11pt]{article}
%
% Documents must define before \input:
%   \newcommand{\doctitle}{...}       Full document title
%   \newcommand{\shorttitle}{...}     Short title for header
%   \newcommand{\docdate}{YYYY-MM-DD} Document date

\usepackage[margin=1in]{geometry}
\usepackage{amsmath,amssymb}
\usepackage[T1]{fontenc}
\usepackage{lmodern}
\usepackage{microtype}
\usepackage{graphicx}
\graphicspath{{figures/}}
\usepackage{booktabs}
\usepackage{longtable}
\usepackage{array}
\usepackage{hyperref}
\usepackage{xcolor}
\usepackage{fancyhdr}
\usepackage{parskip}
\usepackage{caption}
\usepackage{enumitem}
\usepackage{verbatim}
\usepackage{listings}

\lstset{
  language=R,
  basicstyle=\small\ttfamily,
  keywordstyle=\color{blue!70!black},
  commentstyle=\color{green!50!black},
  stringstyle=\color{red!60!black},
  breaklines=true,
  frame=single,
  framerule=0.5pt,
  rulecolor=\color{gray!40},
  backgroundcolor=\color{gray!5},
  xleftmargin=1em,
  framexleftmargin=1em,
  columns=flexible
}

\hypersetup{
  colorlinks=true,
  linkcolor=blue!60!black,
  citecolor=green!50!black,
  urlcolor=blue!60!black
}

\pagestyle{fancy}
\fancyhf{}
\lhead{\footnotesize pmsimstats-ng Technical Report}
\rhead{\footnotesize \shorttitle}
\rfoot{\footnotesize \thepage}
\renewcommand{\headrulewidth}{0.4pt}

\title{\doctitle}
\author{pmsimstats team}
\date{\docdate}
 after \documentclass[11pt]{article}
%
% Documents must define before \input:
%   \newcommand{\doctitle}{...}       Full document title
%   \newcommand{\shorttitle}{...}     Short title for header
%   \newcommand{\docdate}{YYYY-MM-DD} Document date

\usepackage[margin=1in]{geometry}
\usepackage{amsmath,amssymb}
\usepackage[T1]{fontenc}
\usepackage{lmodern}
\usepackage{microtype}
\usepackage{graphicx}
\graphicspath{{figures/}}
\usepackage{booktabs}
\usepackage{longtable}
\usepackage{array}
\usepackage{hyperref}
\usepackage{xcolor}
\usepackage{fancyhdr}
\usepackage{parskip}
\usepackage{caption}
\usepackage{enumitem}
\usepackage{verbatim}
\usepackage{listings}

\lstset{
  language=R,
  basicstyle=\small\ttfamily,
  keywordstyle=\color{blue!70!black},
  commentstyle=\color{green!50!black},
  stringstyle=\color{red!60!black},
  breaklines=true,
  frame=single,
  framerule=0.5pt,
  rulecolor=\color{gray!40},
  backgroundcolor=\color{gray!5},
  xleftmargin=1em,
  framexleftmargin=1em,
  columns=flexible
}

\hypersetup{
  colorlinks=true,
  linkcolor=blue!60!black,
  citecolor=green!50!black,
  urlcolor=blue!60!black
}

\pagestyle{fancy}
\fancyhf{}
\lhead{\footnotesize pmsimstats-ng Technical Report}
\rhead{\footnotesize \shorttitle}
\rfoot{\footnotesize \thepage}
\renewcommand{\headrulewidth}{0.4pt}

\title{\doctitle}
\author{pmsimstats team}
\date{\docdate}
 after \documentclass[11pt]{article}
%
% Documents must define before \input:
%   \newcommand{\doctitle}{...}       Full document title
%   \newcommand{\shorttitle}{...}     Short title for header
%   \newcommand{\docdate}{YYYY-MM-DD} Document date

\usepackage[margin=1in]{geometry}
\usepackage{amsmath,amssymb}
\usepackage[T1]{fontenc}
\usepackage{lmodern}
\usepackage{microtype}
\usepackage{graphicx}
\graphicspath{{figures/}}
\usepackage{booktabs}
\usepackage{longtable}
\usepackage{array}
\usepackage{hyperref}
\usepackage{xcolor}
\usepackage{fancyhdr}
\usepackage{parskip}
\usepackage{caption}
\usepackage{enumitem}
\usepackage{verbatim}
\usepackage{listings}

\lstset{
  language=R,
  basicstyle=\small\ttfamily,
  keywordstyle=\color{blue!70!black},
  commentstyle=\color{green!50!black},
  stringstyle=\color{red!60!black},
  breaklines=true,
  frame=single,
  framerule=0.5pt,
  rulecolor=\color{gray!40},
  backgroundcolor=\color{gray!5},
  xleftmargin=1em,
  framexleftmargin=1em,
  columns=flexible
}

\hypersetup{
  colorlinks=true,
  linkcolor=blue!60!black,
  citecolor=green!50!black,
  urlcolor=blue!60!black
}

\pagestyle{fancy}
\fancyhf{}
\lhead{\footnotesize pmsimstats-ng Technical Report}
\rhead{\footnotesize \shorttitle}
\rfoot{\footnotesize \thepage}
\renewcommand{\headrulewidth}{0.4pt}

\title{\doctitle}
\author{pmsimstats team}
\date{\docdate}


\begin{document}
\maketitle
\thispagestyle{fancy}

\begin{abstract}
When the \texttt{pmsimstats} data-generating process was
updated from compound symmetry to AR(1) within-factor
autocorrelation to expand the feasible parameter space, the
analysis model's assumption of independent residuals (given
a random intercept) became misspecified. This misspecification
inflated Type~I error rates for the OL and CO designs from the
nominal 5\% to 13--17\%. This paper documents the problem,
derives the appropriate analysis model, and validates the fix
by switching from \texttt{lme4::lmer} to \texttt{nlme::lme}
with a continuous-time AR(1) residual correlation structure
(\texttt{corCAR1}).
\end{abstract}

\bigskip
\noindent\fbox{\parbox{\dimexpr\textwidth-2\fboxsep-2\fboxrule}{%
\textbf{Architecture scope.} The AR(1) autocorrelation structure
described here applies to both DGP architectures (A and~B), as it
governs the within-factor temporal correlation in the covariance
matrix regardless of whether the biomarker-treatment interaction is
in the covariance or mean structure. The Type~I error analysis was
conducted under Architecture~B (MVN differential correlation). For
a comparison of both DGP architectures, see
\texttt{02-dgp-mean-moderation-vs-mvn.tex}. For implementation
details of Architecture~A, see
\texttt{19-mean-moderation-implementation-notes.tex}.}}
\bigskip

\tableofcontents
\newpage

% =================================================================
\section{The Problem}
% =================================================================

\subsection{Background}

The original \texttt{pmsimstats} DGP used compound symmetry
(CS) autocorrelation: the correlation between the same factor
at any two timepoints was a constant $\rho = 0.8$, regardless
of the time gap between them. The analysis model used
\texttt{lme4::lmer} with a random intercept per participant
and no explicit residual correlation structure:

\begin{verbatim}
Sx ~ bm + t + Dbc + bm:Dbc + (1|ptID)     # varInDb
Sx ~ bm + t + bm:t + (1|ptID)             # OL (no varInDb)
\end{verbatim}

Under CS, the random intercept absorbs the constant
within-person correlation, and the residuals are approximately
independent. Type~I error rates were near the nominal 5\%.

\subsection{The AR(1) Change}

To expand the feasible biomarker correlation range (from
$c_{bm} \leq 0.25$ to $c_{bm} \leq 0.5$), the DGP was
updated to use AR(1) autocorrelation:

\[
\text{Cor}(\text{Factor}_{t_i}, \text{Factor}_{t_j}) =
\rho^{|t_i - t_j|}
\]

where $t$ is measured in weeks. At $\rho = 0.7$, adjacent
timepoints 2.5 weeks apart have correlation
$0.7^{2.5} = 0.41$, while timepoints 20 weeks apart have
correlation $0.7^{20} = 0.0008$.

Unlike CS, AR(1) produces time-varying correlations that the
random intercept cannot fully absorb. The residuals after
removing the random intercept retain autocorrelation that
decays with time, and the \texttt{lmer} model (which assumes
conditionally independent residuals) underestimates the
standard errors of the fixed effects.

\subsection{Observed Type I Error Inflation}

With the AR(1) DGP and the original \texttt{lmer} analysis
model, Type~I error rates at $c_{bm} = 0$ were:

\begin{table}[h]
\centering
\caption{Type~I error rates with AR(1) DGP and \texttt{lmer}
  analysis (200 replicates, no censoring)}
\begin{tabular}{lrr}
\toprule
Design & $N = 35$ & $N = 70$ \\
\midrule
OL & 0.13 & 0.09 \\
OL+BDC & 0.09 & 0.05 \\
CO & \textbf{0.17} & \textbf{0.13} \\
N-of-1 & 0.06 & 0.08 \\
\bottomrule
\end{tabular}
\end{table}

The CO design showed 17\% false positive rate at $N = 35$
(3.4$\times$ the nominal rate). OL was similarly inflated.
OL+BDC and N-of-1 were closer to nominal.

\subsection{Why CO and OL Are Most Affected}

\textbf{CO design.} The crossover has two paths: drug-first
(Path A: on-on-on-on-off-off-off-off) and placebo-first
(Path B: off-off-off-off-on-on-on-on). Under AR(1), the
autocorrelation within each path decays with time, creating
a correlation gradient across timepoints. The \texttt{Dbc}
variable (binary on/off drug) is perfectly confounded with
the first vs.\ second half of the trial. The \texttt{bm:Dbc}
interaction can absorb spurious signal from the time-dependent
correlation structure that the random intercept does not
capture.

\textbf{OL design.} All participants are always on drug, so
the model uses \texttt{bm:t} (biomarker $\times$ time) as
the interaction term. Under AR(1), the residual correlation
decays with time, creating time-dependent heteroscedasticity
in the residuals that inflates the \texttt{bm:t} standard
error estimate downward.

\textbf{OL+BDC and N-of-1.} These designs have more complex
path structures with multiple transitions between on-drug and
off-drug. The \texttt{Dbc} variable is less confounded with
linear time, so the AR(1) residual structure has less impact
on the interaction estimate.

% =================================================================
\section{The Solution: \texttt{nlme::lme} with \texttt{corCAR1}}
% =================================================================

\subsection{Continuous-Time AR(1)}

The \texttt{nlme} package provides the \texttt{corCAR1}
correlation structure, which models continuous-time AR(1)
residual correlation:

\[
\text{Cor}(\varepsilon_{i,t_1}, \varepsilon_{i,t_2}) =
\phi^{|t_1 - t_2|}
\]

where $\phi$ is estimated from the data. This directly matches
the AR(1) structure used in the DGP (with $\phi$ corresponding
to the DGP's $\rho$).

\subsection{Model Specification}

The analysis model was changed from \texttt{lme4::lmer}:

\begin{verbatim}
lmer(Sx ~ bm + t + Dbc + bm:Dbc + (1|ptID),
     data = datamerged)
\end{verbatim}

to \texttt{nlme::lme}:

\begin{verbatim}
lme(Sx ~ bm + t + Dbc + bm:Dbc,
    random = ~1|ptID,
    correlation = corCAR1(form = ~t|ptID),
    data = datamerged)
\end{verbatim}

The key addition is
\texttt{correlation = corCAR1(form = \textasciitilde t|ptID)},
which tells \texttt{nlme} to estimate and account for AR(1)
residual correlation as a function of the time variable
\texttt{t}, grouped by participant \texttt{ptID}.

\subsection{Coefficient Extraction}

The \texttt{nlme::lme} summary uses a different output
format than \texttt{lme4::lmer}:

\begin{itemize}
\item \texttt{lmer}: \texttt{summary(fit)\$coefficients}
  with columns \texttt{Estimate}, \texttt{Std. Error},
  \texttt{Pr(>|t|)}
\item \texttt{lme}: \texttt{summary(fit)\$tTable} with
  columns \texttt{Value}, \texttt{Std.Error},
  \texttt{p-value}
\end{itemize}

The code was updated to extract from the \texttt{tTable}
format.

\subsection{Error Handling}

Two levels of fallback were implemented:

\begin{enumerate}
\item If \texttt{lme} with \texttt{corCAR1} fails to
  converge (can happen with small sample sizes or
  degenerate data), the model falls back to \texttt{lme}
  without a correlation structure.
\item If the fallback also fails, the function returns
  \texttt{NA} for all output values.
\end{enumerate}

Additionally, rows with \texttt{NA} outcomes (from censoring)
are removed before fitting, as \texttt{nlme::lme} uses
\texttt{na.fail} by default.

% =================================================================
\section{Why This Is the Correct Model}
% =================================================================

\subsection{DGP-Analysis Consistency}

A simulation study's power estimates are valid when the
analysis model is correctly specified relative to the DGP.
When the DGP generates data with AR(1) residual
autocorrelation, the analysis model must account for this
structure. Failure to do so produces:

\begin{itemize}
\item \textbf{Underestimated standard errors:} The effective
  sample size is smaller than the nominal sample size because
  correlated observations carry less independent information.
  Ignoring the correlation inflates the apparent precision.
\item \textbf{Inflated Type~I error:} The underestimated
  standard errors produce test statistics that are too large,
  leading to excess false positives.
\item \textbf{Biased power estimates:} Power appears higher
  than it actually is because the test rejects too often
  under both the null and alternative hypotheses.
\end{itemize}

\subsection{The Random Intercept Is Not Sufficient}

Under compound symmetry, the within-person correlation is
constant across all timepoint pairs. A random intercept
$u_i \sim N(0, \sigma_u^2)$ induces an equicorrelation
structure:

\[
\text{Cor}(Y_{i,t_1}, Y_{i,t_2}) =
\frac{\sigma_u^2}{\sigma_u^2 + \sigma_\varepsilon^2}
= \text{constant}
\]

This exactly matches CS, so the residuals after removing the
random effect are independent.

Under AR(1), the within-person correlation varies:
$\rho^{|t_1 - t_2|}$. The random intercept can absorb the
``average'' correlation but leaves a residual correlation
pattern:

\[
\text{Cor}(\varepsilon_{i,t_1}, \varepsilon_{i,t_2} \mid u_i) =
\rho^{|t_1 - t_2|} - \frac{\sigma_u^2}{\sigma_u^2 +
\sigma_\varepsilon^2}
\]

This residual correlation is positive for nearby timepoints
(where $\rho^{|t_1 - t_2|}$ exceeds the average) and negative
for distant timepoints (where it falls below). The
\texttt{lmer} model treats these residuals as independent,
producing incorrect standard errors.

\subsection{Why \texttt{corCAR1} Is Preferred Over Discrete
            AR(1)}

The timepoints in the trial designs are not equally spaced.
For example, the N-of-1 design has timepoints at weeks
4, 8, 9, 10, 11, 12, 16, 20 --- with gaps ranging from
1 to 4 weeks. The continuous-time AR(1) model
(\texttt{corCAR1}) handles unequal spacing naturally:
$\phi^{|t_1 - t_2|}$ where $t$ is in weeks. The discrete
AR(1) (\texttt{corAR1}) would assume equal spacing between
consecutive observations, which is incorrect for these
designs.

% =================================================================
\section{Design-Specific Model Selection}
% =================================================================

The analysis model varies by design based on the structure
of the drug exposure variable:

\subsection{Designs with Within-Subject Drug Variation
            (CO, N-of-1, OL+BDC)}

These designs have participants who experience both on-drug
and off-drug timepoints. The model tests whether the
biomarker predicts differential response to drug vs.\
no-drug:

\begin{equation}
S_{i,t} = \beta_0 + \beta_1 \cdot bm_i + \beta_2 \cdot t +
\beta_3 \cdot D_{bc,t} +
\beta_4 \cdot bm_i \cdot D_{bc,t} + u_i +
\varepsilon_{i,t}
\label{eq:varInDb}
\end{equation}

where $D_{bc}$ is the binary drug indicator and $\beta_4$ is
the biomarker--treatment interaction of interest.
$\varepsilon_{i,t}$ follows continuous-time AR(1) with
parameter $\phi$ estimated from the data.

\subsection{Open-Label Design (OL)}

All participants are always on drug, so there is no
within-subject drug variation. The model tests whether the
biomarker predicts the rate of improvement over time:

\begin{equation}
S_{i,t} = \beta_0 + \beta_1 \cdot bm_i + \beta_2 \cdot t +
\beta_3 \cdot bm_i \cdot t + u_i + \varepsilon_{i,t}
\label{eq:OL}
\end{equation}

where $\beta_3$ is the biomarker--time interaction. The AR(1)
correlation structure is particularly important here because
\texttt{bm:t} is directly confounded with the time-dependent
residual correlation under AR(1).

% =================================================================
\section{Implementation}
% =================================================================

The complete model fitting code:

\begin{verbatim}
# Build fixed-effects formula
modelbase <- "Sx ~ bm + t"
if (varInDb) {
  modelbase <- paste0(modelbase, " + Dbc + bm:Dbc")
} else {
  modelbase <- paste0(modelbase, " + bm:t")
}
form <- as.formula(modelbase)

# Fit with nlme::lme and AR(1) residual correlation
fit <- tryCatch(
  nlme::lme(form, random = ~1|ptID,
    correlation = nlme::corCAR1(form = ~t|ptID),
    data = datamerged,
    control = nlme::lmeControl(
      opt = "optim",
      maxIter = 200, msMaxIter = 200)),
  error = function(e) {
    # Fallback: no correlation structure
    tryCatch(
      nlme::lme(form, random = ~1|ptID,
        data = datamerged,
        control = nlme::lmeControl(opt = "optim")),
      error = function(e2) NULL
    )
  }
)

# Extract from nlme tTable
c <- summary(fit)$tTable
p <- c[target, "p-value"]
beta <- c[target, "Value"]
betaSE <- c[target, "Std.Error"]
\end{verbatim}

% =================================================================
\section{Validation}
% =================================================================

The 200-replicate simulation confirms that the \texttt{corCAR1}
model corrects the Type~I error inflation:

\begin{table}[h]
\centering
\caption{Type~I error rates at $c_{bm} = 0$, no censoring,
  carryover $= 0$ (200 replicates). The \texttt{corCAR1}
  model restores rates to the nominal 5\% range.}
\begin{tabular}{lrrrr}
\toprule
& \multicolumn{2}{c}{\texttt{lmer} (before)}
& \multicolumn{2}{c}{\texttt{lme + corCAR1} (after)} \\
\cmidrule(lr){2-3} \cmidrule(lr){4-5}
Design & $N = 35$ & $N = 70$ & $N = 35$ & $N = 70$ \\
\midrule
OL & 0.13 & 0.09 & \textbf{0.03} & \textbf{0.02} \\
OL+BDC & 0.09 & 0.05 & \textbf{0.03} & \textbf{0.04} \\
CO & 0.17 & 0.13 & \textbf{0.06} & \textbf{0.08} \\
N-of-1 & 0.06 & 0.08 & \textbf{0.04} & \textbf{0.06} \\
\bottomrule
\end{tabular}
\end{table}

All designs are now within the expected range for 200
replicates (a 95\% binomial confidence interval for
$p = 0.05$ at $n = 200$ is $[0.02, 0.09]$).

Power estimates with carryover also show the expected
pattern:

\begin{table}[h]
\centering
\caption{Power at $c_{bm} = 0.5$, no censoring, with
  \texttt{lme + corCAR1} analysis model (200 replicates).
  $\lambda_{cor} = \ln 2 / t_{1/2}$ (auto).%
  \footnote{Note: $c_{bm} = 0.5$ slightly exceeds the
  PD-feasible ceiling of approximately 0.45 for some trial
  designs under Architecture~B. Results at this value may
  reflect silent covariance matrix correction.}}
\begin{tabular}{llrrrl}
\toprule
Design & $N$ & $t_{1/2} = 0$ & $t_{1/2} = 0.5$ &
  $t_{1/2} = 1.0$ & Carryover effect \\
\midrule
OL & 35 & 0.08 & 0.06 & 0.09 & None (always on drug) \\
OL & 70 & 0.12 & 0.14 & 0.15 & None \\
\midrule
OL+BDC & 35 & 0.34 & 0.24 & 0.18 & Moderate drop \\
OL+BDC & 70 & 0.68 & 0.55 & 0.27 & \textbf{Large drop} \\
\midrule
CO & 35 & 0.27 & 0.25 & 0.19 & Small drop \\
CO & 70 & 0.43 & 0.44 & 0.34 & Small drop \\
\midrule
N-of-1 & 35 & 0.53 & 0.33 & 0.23 & \textbf{Large drop} \\
N-of-1 & 70 & 0.83 & 0.74 & 0.55 & \textbf{Large drop} \\
\bottomrule
\end{tabular}
\end{table}

The paper's central finding is confirmed with correct Type~I
error control: N-of-1 and OL+BDC designs are sensitive to
carryover (power drops of 0.28--0.41 at $t_{1/2} = 1.0$
week), while the CO design is more robust (drops of
0.08--0.10). The power drops are moderate and clinically
interpretable at realistic half-lives, unlike the
artifact-driven crash (0.95 $\to$ 0.18) in the original
publication at $t_{1/2} = 0.1$ weeks.

% =================================================================
\section{Summary of Changes Across the Analysis Model History}
% =================================================================

\begin{table}[h]
\centering
\caption{Evolution of the analysis model}
\small
\begin{tabular}{llll}
\toprule
Version & Package & Correlation & Issue \\
\midrule
Original (42ac030) & \texttt{lme4::lmer} & None (RI only) &
  Correct under CS \\
Ron Thomas (8609f12) & \texttt{lme4::lmer} & None (RI only) &
  Correct under CS \\
HEAD (pre-fix) & \texttt{lme4::lmer} & None (RI only) &
  Inflated Type~I under AR(1) \\
HEAD (current) & \texttt{nlme::lme} & \texttt{corCAR1} &
  Correct under AR(1) \\
\bottomrule
\end{tabular}
\end{table}

The transition from \texttt{lmer} to \texttt{lme} is driven
entirely by the need to specify an explicit residual
correlation structure, which \texttt{lme4} does not support.
The fixed-effects formula is unchanged. The random effects
structure (random intercept per participant) is unchanged.
The only addition is the \texttt{corCAR1} term that accounts
for the time-dependent residual autocorrelation introduced by
the AR(1) DGP.

\vfill
\noindent\rule{\textwidth}{0.4pt}
{\footnotesize
Rendered on 2026-03-25 at 18:49 PDT.\\
Source: \textasciitilde/prj/alz/10-pmsimstats-ng/pmsimstats-ng/docs/06-ar1-residual-correlation.tex}

\end{document}
